\section{Related Work}

\noindent\textbf{Text-Guided Video Diffusion Models.}
Generating videos from text instructions using Diffusion Models \cite{DDPM_paper,sohl2015diffusion_theory} is a newly established and actively researched field introduced by Video Diffusion Models (VDM) \cite{video-diffusion-models}. The method can generate only low-resolution videos (up to 128x128) with a maximum of 16 frames (without autoregression), imposing significant limitations, while requiring massive training resources.
%Training text-to-video models usually involves large datasets \cite{Bain21WebVid,wang2023internvid,ju2024miradatalargescalevideodataset,chen2024panda70m}.
Several methods thus employ spatial/temporal upsampling \cite{imagen_video,make-a-video,video-diffusion-models,AlignYourLatents}, using cascades with up to 7 enhancer modules  \cite{imagen_video}. While this leads to high-resolution and long videos, the generated content is still limited by the content depicted in the key frames. 

Towards  generating longer videos (\ie more keyframes), Text-To-Video-Zero (T2V0) \cite{text2video-zero} and ART-V \cite{park2021nerfies_art_v} utilize a text-to-image diffusion model. Thus, they can generate only simple motions. 
T2V0 conditions on its first frame via cross-frame attention and  ART-V on an anchor frame. The lack of global reasoning leads to  unnatural or repetitive motions.  
MTVG \cite{oh2023mtvg} and FIFO-Diffusion \cite{kim2024fifo} transforms a text-to-video model into an autoregressive method through a training-free approach. As it uses strong consistency priors within and between chunks, it results in low motion amount, and mostly near-static background. 
%
%
%It employs strong consistency priors between and among video chunks, which leads to very low motion amount, and mostly near-static background.  
FreeNoise \cite{qiu2023freenoise} samples a small set of noise vectors, re-uses them for the generation of all frames, while temporal attention is performed on local windows. As temporal attention is invariant to such frame shuffling, it leads to high similarity between frames, almost always static global motion and near-constant videos.  
Gen-L \cite{wang2023gen_L} generates overlapping short videos and combines them via temporal co-denoising, which can lead to quality degradation with video stagnation. Recent transformed-based diffusion models  \cite{opensora,opensoraplan} operate
in the latent space of a 3D autoencoder, enabling the generation of up to 384 frames. Despite extensive training,
these models produce videos with limited motion, often resulting in near-constant videos.
%This suggests that further significant training effort is need when utilizing the transformer architecture.

\noindent\textbf{Image-Guided Video Diffusion Models as Long Video Generators.}
Several works condition the video generation by a driving image or video  \cite{xing2023dynamicrafter,chen2023seine,SVD,guo2023sparsectrl,gen1_paper,2023i2vgenxl,chen2023livephoto,MakePixelsDance,dai2023animateanything,wang2024videocomposer,long2024videodrafter,ren2024consisti2v}. They can thus be turned into an autoregressive method by conditioning on the frame(s) of the preceding chunk.


%Some video diffusion models are conditioned on key frames (\eg generated by a text-to-image model) to independently generate multiple videos with shared context.

VideoDrafter \cite{long2024videodrafter} takes an anchor frames (from a  text-to-image model) and conditions  a video diffusion model on it to generate independently multiple videos that share the same high-level context. However, this leads to drastic scene cuts as no consistency among video chunks is enforced. StoryDiffusion \cite{Zhou2024storydiffusion} conditions on video frames that have been linearly propagated from key frames, which leads to severe quality degradation.
Several works \cite{chen2023seine,MakePixelsDance,dai2023animateanything} concatenate the (encoded) conditionings (\eg  input frame(s)) with an additional mask (indicating the provided frame(s)) to the input of the video diffusion model. 

\begin{figure*}[t]
    \centering
    \includegraphics[width=1.0\linewidth]{figures/StreamingT2V.png}
    \caption{
    The overall pipeline of StreamingT2V involves three stages: (i) \textit{Initialization Stage}: The first 16-frame chunk is synthesized by an off-the-shelf text-to-video model. (ii) \textit{Streaming T2V Stage}: New content for subsequent frames is autoregressively generated. (iii) \textit{Streaming Refinement Stage}: The long video (\eg 240, 1200 frames or more) is autoregressively enhanced using a high-resolution text-to-short-video model with a randomized blending approach.
%2. Streaming T2V Stage: New content for subsequent frames is autoregressively generated.
%3. Streaming Refinement Stage: The long video (\eg 240, 1200 frames or more) is autoregressively enhanced using a high-resolution text-to-short-video model (\eg MS-Vid2Vid-XL) with a randomized blending approach.
        %The overall pipeline of StreamingT2V: In the \textit{Initialization Stage} the first 16-frame chunk is synthesized by a text-to-video model (\eg Modelscope \cite{Modelscope}). In the \textit{Streaming T2V Stage} the new content for further frames are autoregressively generated. Finally, in the \textit{Streaming Refinement Stage} the generated long video (\eg 240, 1200 or more frames) is autoregressively enhanced by applying a high-resolution text-to-short-video model (\eg MS-Vid2Vid-XL \cite{2023i2vgenxl}) equipped with our randomized blending approach.
    }
    \label{fig:high_level_pipeline}
\end{figure*}


In addition to concatenating the conditioning to the input of the diffusion model, several works  \cite{SVD,2023i2vgenxl,wang2024videocomposer} replace the text embeddings in the cross-attentions of the diffusion model by CLIP  \cite{CLIP_paper} image embeddings of the conditional frames.  However, according to our experiments,  their applicability for long video generation is limited. SVD \cite{SVD} shows severe quality degradation over time (see \figrefcustom{fig:comparison-to-sota}), and both, I2VGen-XL \cite{2023i2vgenxl} and  SVD \cite{SVD}  generate often inconsistencies between chunks, still indicating that the conditioning mechanism is too weak. 

Some works \cite{xing2023dynamicrafter,chen2023livephoto} such as DynamiCrafter-XL \cite{xing2023dynamicrafter} thus add to each text cross-attention an image cross-attention, which leads to better quality, but still to frequent inconsistencies between chunks.

The concurrent work SparseCtrl \cite{guo2023sparsectrl} adds a ControlNet \cite{controlnet} like branch to the model, using as input the conditional frames and a frame-mask. By design, it requires to append additional black frames to the conditional frames.  This inconsistency is difficult to compensate for the model, leading to frequent and severe scene cuts between frames.  

% Overall,  only a small number of keyframes can currently be generated at once with high quality. In-between frames can be interpolated, but it does not lead to new content. Image-to-video methods can be used autoregressively, but their conditional mechanisms lead either to inconsistencies, or the method suffers from video stagnation. 
% We conclude that existing works are not suitable for high-quality and consistent long video generation without video stagnation.